% !TEX program = xelatex
% =====================================================================
%  KIẾN TRÚC HỆ THỐNG
%  "Hệ thống ghép cặp CV - Việc làm theo section"
%  (Section-Aware CV-Job Matching + hàng đợi xử lý file nền)
%
%  Biên dịch:  xelatex kien_truc_he_thong.tex      (MiKTeX / TeX Live)
%  Lưu ý: bắt buộc XeLaTeX (fontspec) để hiển thị đúng tiếng Việt.
%
%  Sơ đồ:   8 hình TikZ đã được tách ra thư mục so_do/ (mỗi hình một file
%           .tex, file này chỉ \input{so_do/figXX_...}); caption/label nằm
%           cùng file hình nên mọi \ref vẫn hoạt động.
%  Ảnh PNG: python so_do/build_png.py  ->  so_do/png/figXX_....png (200 dpi)
% =====================================================================
\documentclass[a4paper,11pt]{article}

\usepackage[a4paper,top=2cm,bottom=2cm,left=2cm,right=2cm]{geometry}
% Tiếng Việt không có dấu ngắt từ tự nhiên trong XeLaTeX -> nới giãn một chút
% để tránh các dòng tràn ra ngoài lề.
\setlength{\emergencystretch}{3em}
\usepackage{fontspec}
\setmainfont{Times New Roman}
\setsansfont{Arial}
\setmonofont[Scale=0.85]{Consolas}
\usepackage{amsmath,amssymb}
\usepackage{booktabs}
\usepackage{array}
\usepackage{enumitem}
\usepackage{xcolor}
\usepackage{caption}
\captionsetup{font=small,labelfont=bf}
\usepackage{tikz}
\usetikzlibrary{arrows.meta,positioning,fit,calc,automata,shapes.geometric}
\usepackage[colorlinks=true,linkcolor=blue!50!black,urlcolor=blue!70!black,
            citecolor=green!40!black]{hyperref}

% ---------------- màu + style dùng chung cho sơ đồ ----------------
\definecolor{cblue}{HTML}{2B6CB0}
\definecolor{cbluel}{HTML}{EAF2FB}
\definecolor{cgreen}{HTML}{276749}
\definecolor{cgreenl}{HTML}{E6F4EC}
\definecolor{corange}{HTML}{B7791F}
\definecolor{corangel}{HTML}{FBF3E0}
\definecolor{cpurple}{HTML}{6B46C1}
\definecolor{cpurplel}{HTML}{F0E9FA}
\definecolor{cred}{HTML}{C53030}
\definecolor{credl}{HTML}{FDE8E8}
\definecolor{cgray}{HTML}{4A5568}
\definecolor{cgrayl}{HTML}{EDF2F7}

\tikzset{
  svc/.style   = {draw=cblue, line width=0.7pt, rounded corners=3pt, fill=cbluel,
                  inner sep=4pt, font=\scriptsize, align=center},
  infra/.style = {draw=cgreen, line width=0.7pt, rounded corners=3pt, fill=cgreenl,
                  inner sep=4pt, font=\scriptsize, align=center},
  store/.style = {draw=corange, line width=0.7pt, rounded corners=3pt, fill=corangel,
                  inner sep=4pt, font=\scriptsize, align=center},
  llm/.style   = {draw=cpurple, line width=0.7pt, rounded corners=3pt, fill=cpurplel,
                  inner sep=4pt, font=\scriptsize, align=center},
  fail/.style  = {draw=cred, line width=0.7pt, rounded corners=3pt, fill=credl,
                  inner sep=4pt, font=\scriptsize, align=center},
  box/.style   = {draw=black!45, rounded corners=2pt, fill=white, inner sep=3.5pt,
                  font=\scriptsize, align=center},
  band/.style  = {draw=black!35, rounded corners=4pt, fill=black!4, inner sep=6pt,
                  align=left},
  arr/.style   = {-{Stealth[length=2.2mm]}, semithick},
  arrl/.style  = {-{Stealth[length=2.4mm]}, thick},
  arrd/.style  = {-{Stealth[length=2.2mm]}, semithick, dashed, black!60},
  lbl/.style   = {font=\tiny, align=center, inner sep=1.5pt},
  actor/.style = {draw=black!55, fill=cgrayl, rounded corners=2pt, minimum width=2.1cm,
                  minimum height=0.7cm, font=\scriptsize, align=center},
  st/.style    = {state, minimum size=1.2cm, font=\scriptsize, align=center, inner sep=1pt},
}

\title{\textbf{KIẾN TRÚC HỆ THỐNG}\\[4pt]
       \large Hệ thống ghép cặp CV -- Việc làm theo section\\
       \normalsize (Section-Aware CV--Job Matching) kèm hàng đợi xử lý file nền}
\author{Nhóm nghiên cứu: .........................\\
        \small Giáo viên hướng dẫn: ........................}
\date{\today}

\begin{document}
\maketitle

\begin{abstract}
\noindent Tài liệu mô tả kiến trúc phần mềm của hệ thống ghép cặp ``sơ yếu lý lịch (CV) --
việc làm (JD)'' theo từng \emph{phần (section)} của văn bản, gồm hai tiến trình lớn:
(1) \textbf{AI service} viết bằng Python/FastAPI, đảm nhận trích xuất văn bản, phân tích
bằng mô hình ngôn ngữ lớn (LLM), sinh embedding theo section, lưu vector và xếp hạng
hybrid (semantic + từ khóa + ràng buộc có cấu trúc); và (2) \textbf{API chính} viết bằng
NestJS, đảm nhận nghiệp vụ người dùng và \emph{gọi} AI service. Trọng tâm của tài liệu là
\textbf{hàng đợi xử lý file nền} (background processing queue) cho phép người dùng upload
CV và nhận trạng thái xử lý từng bước thay cho lời gọi đồng bộ ``treo'' hàng chục giây.
Tài liệu gồm: sơ đồ tổng thể, sơ đồ lớp module, máy trạng thái job, sơ đồ chuỗi lời gọi
(bước upload bất đồng bộ), mô hình dữ liệu vector theo section, và \textbf{luồng ghép cặp
(matching) trình bày theo dạng dễ hiểu}: bốn bước từ tìm ứng viên đến trả kết quả, cách
tính điểm mức độ hợp, và hai chiều ghép cặp \texttt{jd2cv}/\texttt{cv2jd}; phần cuối nêu
hướng nghiên cứu khoa học đề xuất.
\end{abstract}

\tableofcontents

% =====================================================================
\section{Tổng quan hệ thống}
% =====================================================================

Hệ thống có hai tiến trình chính (hai mã nguồn khác nhau) và một số thành phần hạ tầng:

\begin{itemize}[leftmargin=1.2em]
  \item \textbf{API chính (NestJS, TypeScript)} --- backend nghiệp vụ: xác thực, hồ sơ
        người dùng, tin tuyển dụng, hợp đồng, thanh toán. Là nơi duy nhất được phép
        đọc/ghi \textbf{PostgreSQL} (nguồn dữ liệu gốc).
  \item \textbf{AI service (Python, FastAPI)} --- microservice nội bộ, \emph{không} có
        quyền ghi PostgreSQL: trích xuất PDF/DOCX, phân tích cấu trúc bằng
        \texttt{qwen3:4b} (chạy local qua Ollama), sinh embedding \texttt{BAAI/bge-m3},
        lưu vector vào \textbf{ChromaDB}, xếp hạng và giải thích bằng LLM.
  \item \textbf{ChromaDB} --- cơ sở vector (HNSW, cosine) persistent trên đĩa.
  \item \textbf{SQLite} --- bảng \texttt{processing.db} lưu \emph{trạng thái job xử lý
        file} (hàng đợi nền). Lý do dùng SQLite: ở giai đoạn này AI service chỉ được cấp
        quyền \emph{đọc} PostgreSQL; giao diện \texttt{ProcessingStore} được thiết kế để
        sau này thay bằng \texttt{PostgresProcessingStore} mà không đổi pipeline.
  \item \textbf{Outbox + \texttt{sync\_worker}} --- cơ chế đồng bộ \emph{một chiều}
        PostgreSQL $\rightarrow$ Chroma: các bản ghi CV/JD trong DB được đẩy sang vector
        store để phục vụ truy hồi.
\end{itemize}

\begin{figure}[htbp]
\centering
\resizebox{\textwidth}{!}{%
\begin{tikzpicture}[node distance=0.9cm and 1.1cm, font=\scriptsize]
  % -------- dải 0: ứng dụng --------
  \node[svc]   (web)     at (0,0)     {Web / Mobile\\(Next.js)};
  \node[svc]   (nest)    at (5.5,0)   {API chính\\NestJS + Prisma};
  \node[infra] (pg)      at (16,0)    {PostgreSQL\\(nguồn gốc)};

  % -------- dải 1: xử lý --------
  \node[svc]   (api)     at (5.5,-2.3){AI Service\\FastAPI \texttt{/internal/*}};
  \node[svc]   (worker)  at (7,-4.6)  {Worker nền (1 thread)\\claim $\to$ pipeline $\to$ ghi kết quả};
  \node[store] (sqlite)  at (12,-4.6) {SQLite\\bảng \texttt{jobs}};
  \node[store] (sync)    at (16,-2.2) {\texttt{sync\_worker}\\(đọc outbox)};

  % -------- dải 2: hạ tầng --------
  \node[llm]   (ollama)  at (3.3,-7.4){Ollama\\LLM \texttt{qwen3:4b}};
  \node[llm]   (bge)     at (7.6,-7.4){Embedder\\BGE-M3 (CPU)};
  \node[infra] (chroma)  at (12.2,-7.4){ChromaDB\\vector theo section};
  \node[store] (files)   at (16.2,-7.4){File tạm\\(PDF/DOCX)};

  % -------- luồng chính --------
  \draw[arr] (web)   -- node[lbl,above]{HTTP/REST} (nest);
  \draw[arr] (nest)  -- node[lbl,right]{giao job,\\poll trạng thái,\\recommend} (api);
  \draw[arr,<->] (nest) -- node[lbl,above]{CRUD; outbox} (pg);
  \draw[arr] (pg)   -- node[lbl,right]{\texttt{outbox.jsonl}} (sync);
  \draw[arr] (sync) -- node[lbl,right]{index / xoa vector} (chroma);
  \draw[arr] (api)  -- node[lbl,left]{khoi tao} (worker);
  \draw[arr] (worker) -- node[lbl,above]{claim, checkpoint, READY} (sqlite);
  \draw[arr] (worker) -- node[lbl,left]{parse JSON} (ollama);
  \draw[arr] (worker) -- node[lbl,left]{nhung section} (bge);
  \draw[arr] (worker) -- node[lbl,below,sloped]{upsert vector} (chroma);
  \draw[arr] (worker) -- node[lbl,above,sloped]{luu / xoa} (files);
\end{tikzpicture}}
\caption{Tổng quan hệ thống: API chính (NestJS) điều phối, AI service (FastAPI) xử lý
và xếp hạng, dữ liệu vector nằm ở ChromaDB, trạng thái job nằm ở SQLite.}
\label{fig:tongthe}
\end{figure}


Hai thành phần \emph{giao tiếp} với nhau chỉ qua HTTP nội bộ, được bảo vệ bằng
\texttt{INTERNAL\_API\_KEY} (che dấu \texttt{/internal/*} khỏi Internet). Nhờ đó AI
service có thể thay đổi độc lập (thêm model, thêm queue) mà không ảnh hưởng API chính.

% =====================================================================
\section{Kiến trúc bên trong AI service}
% =====================================================================

AI service được chia thành bốn lớp phụ trách rõ ràng (Hình~\ref{fig:module}), theo
nguyên tắc \emph{API chỉ điều phối, nghiệp vụ nằm ở lớp giữa, lớp dưới lo hạ tầng}.
Thư viện lõi \texttt{src/cvjd} (Python package \texttt{cvjd}) độc lập với web framework
để có thể chạy benchmark từ dòng lệnh.

\begin{figure}[htbp]
\centering
\resizebox{\textwidth}{!}{%
\begin{tikzpicture}[node distance=0.35cm, font=\scriptsize]
  % ---- lop API ----
  \node[band, minimum width=16.4cm] (b1) at (0,0) {};
  \node[font=\bfseries\footnotesize, anchor=west] at ([xshift=0.1cm]b1.west)
        {Lớp 1 -- API \& an ninh (\texttt{apps/ai/app})};
  \foreach \x/\txt in {-6.2/{\texttt{routes.py}\\17 endpoint},
                      -1.9/{\texttt{schemas.py}\\Pydantic v2},
                      2.2/{\texttt{security.py}\\API key, fail-closed},
                      6.4/{\texttt{errors.py}\\mã lỗi + HTTP},
                      10.2/{\texttt{main.py}\\lifespan, CORS, worker}}
     \node[box, fill=cbluel] at (\x,-1.15) {\txt};
  \node[box, fill=cbluel] at (13.6,-1.15) {\texttt{\_\_init\_\_.py}\\cấu hình app};

  % ---- lop nghiep vu ----
  \node[band, minimum width=16.4cm] (b2) at (0,-2.75) {};
  \node[font=\bfseries\footnotesize, anchor=west] at ([xshift=0.1cm]b2.west)
        {Lớp 2 -- Nghiệp vụ};
  \foreach \x/\txt in {-6.4/{\texttt{text.py}\\magic bytes, PDF/DOCX/OCR},
                      -2.5/{\texttt{normalize.py}\\chuẩn hoá CV/JD},
                      1.2/{\texttt{llm.py}\\client Ollama},
                      4.6/{\texttt{embeddings.py}\\BGE-M3},
                      8.4/{\texttt{recommend.py}\\xếp hạng hybrid},
                      12.4/{\texttt{explain.py}\\giải thích + rerank}}
     \node[box, fill=cbluel] at (\x,-3.9) {\txt};

  % ---- lop hang doi & luu tru ----
  \node[band, minimum width=16.4cm] (b3) at (0,-5.5) {};
  \node[font=\bfseries\footnotesize, anchor=west] at ([xshift=0.1cm]b3.west)
        {Lớp 3 -- Hàng đợi xử lý \& lưu trữ};
  \foreach \x/\txt in {-6.4/{\texttt{processing/worker.py}\\thread nền, retry, backoff},
                      -2.5/{\texttt{processing/pipeline.py}\\4 bước EXTRACT$\to$FINALIZE},
                      1.2/{\texttt{processing/store.py}\\SQLite, máy trạng thái},
                      5.4/{\texttt{index\_store.py}\\payload JSONL},
                      9.2/{\texttt{cvjd.storage.chroma}\\upsert / query vector}}
     \node[box, fill=corangel] at (\x,-6.65) {\txt};

  % ---- thu vien loi ----
  \node[band, minimum width=16.4cm] (b4) at (0,-8.3) {};
  \node[font=\bfseries\footnotesize, anchor=west] at ([xshift=0.1cm]b4.west)
        {Lớp 4 -- Thư viện lõi \texttt{src/cvjd}};
  \foreach \x/\txt in {-6.4/{\texttt{schema.py}, \texttt{io.py}\\dạng dữ liệu, JSONL},
                      -2.5/{\texttt{data/}\\payload, preprocessing},
                      1.2/{\texttt{models/}\\BM25, chunking, embedding, matching},
                      5.4/{\texttt{ranking/}\\engine, hybrid, structured, constraints},
                      9.2/{\texttt{retrieval/}, \texttt{storage/}\\retriever, chroma, outbox}}
     \node[box, fill=cgreenl] at (\x,-9.45) {\txt};

  % ---- quan he giua cac lop ----
  \draw[arrd] ([yshift=0.4cm]b1.south) -- ([yshift=-0.4cm]b2.north)
        node[lbl, right, xshift=0.3cm]{gọi hàm nghiệp vụ};
  \draw[arrd] ([yshift=0.4cm]b2.south) -- ([yshift=-0.4cm]b3.north)
        node[lbl, right, xshift=0.3cm]{worker gọi pipeline; recommend gọi Chroma};
  \draw[arrd] ([yshift=0.4cm]b3.south) -- ([yshift=-0.4cm]b4.north)
        node[lbl, right, xshift=0.3cm]{tái sử dụng thư viện lõi};
\end{tikzpicture}}
\caption{Các lớp module của AI service. Lớp ngoài tiếp nhận HTTP, lớp nghiệp vụ xử lý
văn bản/LLM, lớp hàng đợi lo việc chạy nền, thư viện lõi chứa các mô hình và thứ tự
xếp hạng (có thể chạy riêng từ CLI để benchmark).}
\label{fig:module}
\end{figure}


Hai đường xử lý ``đồng bộ'' và ``nền'' dùng chung các khối nghiệp vụ nên kết quả phân
tích là như nhau (Hình~\ref{fig:2duong}).

\begin{figure}[htbp]
\centering
\resizebox{0.92\textwidth}{!}{%
\begin{tikzpicture}[node distance=0.8cm and 0.9cm, font=\scriptsize]
  % --- duong dong bo ---
  \node[box, fill=cbluel, font=\scriptsize\bfseries] (s0) at (0,0)
        {\texttt{POST /api/v1/ai/analyze-cv}\\(đồng bộ, dev/fixture)};
  \node[box, fill=cbluel] (s1) at (5.2,0) {\texttt{text.py}\\đọc file, làm sạch};
  \node[box, fill=cbluel] (s2) at (9.6,0) {\texttt{llm.py} + \texttt{normalize}\\parse JSON (qwen3:4b)};
  \node[box, fill=white]  (s3) at (14.4,0) {JSON trả về\\(không lưu, không index)};

  % --- duong nen ---
  \node[box, fill=corangel, font=\scriptsize\bfseries] (q0) at (0,-2.4)
        {\texttt{POST /internal/jobs}\\(bất đồng bộ, 202)};
  \node[box, fill=corangel] (q1) at (5.2,-2.4) {EXTRACTING\\(worker)};
  \node[box, fill=corangel] (q2) at (9.6,-2.4) {PARSING + EMBEDDING\\(LLM + BGE-M3)};
  \node[box, fill=cgreenl]  (q3) at (14.4,-2.4) {INDEX $\to$ READY\\(Chroma + payload)};

  \draw[arr] (s0) -- (s1); \draw[arr] (s1) -- (s2); \draw[arr] (s2) -- (s3);
  \draw[arr] (q0) -- (q1); \draw[arr] (q1) -- (q2); \draw[arr] (q2) -- (q3);
  \draw[arr] (s1) -- node[lbl,right]{dùng chung} (q1);
  \draw[arr] (s2) -- (q2);
\end{tikzpicture}}
\caption{Hai đường xử lý dùng chung khối trích xuất và LLM: đường đồng bộ
\texttt{analyze-cv} chỉ trả JSON, đường hàng đợi còn nhúng và index vector để phục vụ
truy hồi sau đó.}
\label{fig:2duong}
\end{figure}


% =====================================================================
\section{Hàng đợi xử lý file nền (phần lõi)}
% =====================================================================

\subsection{Vấn đề}

Việc phân tích một CV mất 10--40~giây (trích xuất PDF, gọi LLM sinh JSON, nhúng 5--8
section). Nếu để HTTP đồng bộ ``treo'' chừng đó thời gian thì trình duyệt/người dùng dễ
bị timeout và NestJS cũng bị chiếm kết nối. Giải pháp: \textbf{hàng đợi job nền}
\begin{enumerate}[leftmargin=1.6em]
  \item Client gửi file $\to$ nhận ngay \texttt{202 Accepted} kèm \texttt{job\_id};
  \item Worker chạy nền xử lý từng job, cập nhật \emph{step} + \emph{progress};
  \item Client poll trạng thái và hiển thị ``Đang trích xuất $\to$ Đang phân tích
        $\to$ Đang nhúng $\to$ Hoàn tất''.
\end{enumerate}

\subsection{Máy trạng thái của job}

Mỗi job có \emph{trạng thái} (\texttt{status}) và \emph{bước đang chạy}
(\texttt{step}) là hai chiều độc lập: trạng thái nói ``job đang đi tới đâu về mặt vòng
đời'', bước nói ``đang ở khâu nào'' để hiển thị tiến độ (Hình~\ref{fig:state}).

\begin{figure}[htbp]
\centering
\begin{tikzpicture}[node distance=1.5cm and 1.8cm, font=\scriptsize]
  \node[st, fill=cbluel, draw=cblue]  (p)     at (0,0)    {PENDING};
  \node[st, fill=corangel, draw=corange] (pr)  at (5,0)    {PROCESSING};
  \node[st, fill=cgreenl, draw=cgreen, accepting] (r) at (10,1.7) {READY};
  \node[st, fill=credl, draw=cred]    (f)     at (10,-1.7){FAILED};
  \node[st, fill=cgrayl, draw=cgray]  (c)     at (5,-3.6) {CANCELLED};

  \draw[arr] (p)  -- node[lbl,above]{\texttt{claim()} (attempts $+1$)} (pr);
  \draw[arr] (pr) -- node[lbl,right]{\texttt{mark\_ready()}} (r);
  \draw[arr] (pr) -- node[lbl,right]{lỗi vĩnh viễn\\hoặc hết lượt thử} (f);
  \draw[arr] (pr) to[bend right=25] node[lbl,below]{lỗi tạm thời: backoff $2^{n}\cdot\Delta$} (p);
  \draw[arr] (f)  to[bend right=30] node[lbl,below]{\texttt{POST /jobs/\{id\}/retry}} (p);
  \draw[arr] (p)  -- node[lbl,left]{hủy} (c);
  \draw[arr] (pr) to[bend left=22] node[lbl,above,sloped]{hủy} (c);
  \draw[arr] (pr) edge[loop above, min distance=0.55cm]
        node[lbl,above]{checkpoint \texttt{set\_step()}} (pr);
\end{tikzpicture}
\caption{Máy trạng thái job xử lý file. \texttt{READY} là trạng thái hấp thụ (job hoàn
tất); \texttt{FAILED} vẫn giữ file gốc để có thể retry thủ công.}
\label{fig:state}
\end{figure}


Bên trong \texttt{PROCESSING}, pipeline chạy tuần tự bốn bước, mỗi bước ghi \emph{checkpoint}
vào DB nên nếu service chết giữa chừng thì ta biết chính xác job đang ở đâu:

\begin{figure}[htbp]
\centering
\begin{tikzpicture}[node distance=0.9cm, font=\scriptsize]
  \foreach \i/\nm/\pc in {0/{RECEIVED}/5, 1/{EXTRACTING}/25, 2/{PARSING}/55,
                          3/{EMBEDDING}/85, 4/{FINALIZING}/95} {
    \node[box, fill=cbluel, minimum width=2.4cm] (s\i) at ({\i*3.0},0) {\nm\\[1pt]
          \tiny tiến độ \pc\%};
  }
  \draw[arr] (s0) -- (s1); \draw[arr] (s1) -- (s2);
  \draw[arr] (s2) -- (s3); \draw[arr] (s3) -- (s4);
  \node[lbl, below=0.25cm of s1, text=cgray]{magic bytes, \texttt{extract\_text}, làm sạch};
  \node[lbl, below=0.25cm of s2, text=cgray]{prompt $\to$ JSON (LLM)};
  \node[lbl, below=0.25cm of s3, text=cgray]{nhúng từng section (BGE-M3)};
  \node[lbl, below=0.25cm of s4, text=cgray]{upsert Chroma + payload, \texttt{mark\_ready()}};
\end{tikzpicture}
\caption{Các bước bên trong \texttt{PROCESSING} và mức tiến độ tương ứng được ghi vào
\texttt{step}/\texttt{progress}.}
\label{fig:steps}
\end{figure}


\subsection{Luồng bất đồng bộ khi upload CV}

Hình~\ref{fig:sequence} mô tả toàn bộ lời gọi khi người dùng upload CV: từ HTTP request
đến khi job chuyển sang \texttt{READY} và client poll được kết quả.

\begin{figure}[htbp]
\centering
\resizebox{\textwidth}{!}{%
\begin{tikzpicture}[font=\scriptsize]
  \foreach \i/\nm in {0/{Trình duyệt}, 1/{NestJS API}, 2/{AI Service (HTTP)},
                      3/{Worker (thread)}, 4/{SQLite}, 5/{Ollama + BGE-M3}}{
    \node[actor] (a\i) at ({\i*3.7},0) {\nm};
    \draw[dashed, black!40] (a\i.south) -- ++(0,-11.2);
  }
  \draw[arr] ($(a0)+(0,-0.9)$) -- node[lbl,above]{1. \texttt{POST /api/cvs/upload} (PDF + JD text)} ($(a1)+(0,-0.9)$);
  \draw[arr] ($(a1)+(0,-1.7)$) -- ++(1.35,0) -- node[lbl,right]{2. validate, lưu row CV (đang xử lý)} ++(0,-0.5) -- ++(-1.35,0);
  \draw[arr] ($(a1)+(0,-2.5)$) -- node[lbl,above]{3. \texttt{POST /internal/jobs} (file, documentId, version)} ($(a2)+(0,-2.5)$);
  \draw[arr] ($(a2)+(0,-3.3)$) -- ++(-1.35,0) -- node[lbl,left]{4. \texttt{INSERT} job \texttt{PENDING}; lưu file tạm} ++(0,-0.5) -- ++(1.35,0);
  \draw[arr] ($(a2)+(0,-4.3)$) -- node[lbl,above]{5. \texttt{202 Accepted \{job\_id\}}} ($(a1)+(0,-4.3)$);
  \draw[arr] ($(a3)+(0,-5.1)$) -- node[lbl,above]{6. \texttt{claim}: \texttt{PENDING}$\to$\texttt{PROCESSING}} ($(a4)+(0,-5.1)$);
  \draw[arr] ($(a3)+(0,-5.9)$) -- ++(1.35,0) -- node[lbl,right]{7. EXTRACTING: đọc file, làm sạch} ++(0,-0.5) -- ++(-1.35,0);
  \draw[arr] ($(a3)+(0,-6.9)$) -- node[lbl,above]{8. PARSING: prompt $\to$ JSON (\texttt{qwen3:4b})} ($(a5)+(0,-6.9)$);
  \draw[arr] ($(a3)+(0,-7.7)$) -- node[lbl,above]{9. EMBEDDING: vector từng section} ($(a5)+(0,-7.7)$);
  \draw[arr] ($(a3)+(0,-8.5)$) -- node[lbl,above]{10. upsert vector + payload; \texttt{mark\_ready()}} ($(a4)+(0,-8.5)$);
  \draw[arr] ($(a1)+(0,-9.3)$) -- node[lbl,above]{11. poll \texttt{GET /internal/jobs/\{id\}}} ($(a2)+(0,-9.3)$);
  \draw[arr] ($(a2)+(0,-10.1)$) -- node[lbl,above]{12. \texttt{\{status: READY, result\}}} ($(a1)+(0,-10.1)$);
  \draw[arr] ($(a1)+(0,-10.7)$) -- node[lbl,above]{13. hiển thị CV đã xử lý} ($(a0)+(0,-10.7)$);
\end{tikzpicture}}
\caption{Chuỗi lời gọi của luồng upload bất đồng bộ. Toàn bộ bước 6--10 chạy nền trong
cùng process FastAPI; API chính chỉ poll trạng thái.}
\label{fig:sequence}
\end{figure}


\subsection{Độ tin cậy}

\begin{itemize}[leftmargin=1.2em]
  \item \textbf{Phân loại lỗi:} lỗi \emph{vĩnh viễn} (file hỏng, text rỗng, sai định
        dạng) $\to$ \texttt{FAILED} ngay, không thử lại; lỗi \emph{tạm thời} (Ollama
        timeout, Chroma khóa database) $\to$ worker thử lại với backoff
        $2^{n}\cdot\Delta$ đến \texttt{JOB\_MAX\_ATTEMPTS} (mặc định 3).
  \item \textbf{Phát hiện job ``kẹt'':} nếu \texttt{PROCESSING} quá
        \texttt{JOB\_STALE\_SECONDS} (300~s) mà không cập nhật thì worker tự đưa job về
        \texttt{PENDING}; endpoint \texttt{POST /internal/jobs/reconcile} cho phép vận
        hành viên gỡ kẹt thủ công.
  \item \textbf{Không mất job khi restart:} SQLite là nguồn sự thật duy nhất; khi khởi
        động lại service, các job còn \texttt{PENDING} trong DB sẽ được worker xử lý tiếp.
  \item \textbf{Bỏ qua nhúng lại (idempotent):} nếu nội dung text không đổi
        (\texttt{content\_hash} khớp) \emph{và} vector của phiên bản đó vẫn còn nguyên
        trong Chroma + payload còn trong store thì job chuyển thẳng sang \texttt{READY}
        với cảnh báo \texttt{REUSED\_INDEX} --- retry không tiêu tốn thêm lượt LLM.
  \item \textbf{Quan sát được:} \texttt{GET /internal/jobs/stats} trả về số job theo
        trạng thái, độ sâu hàng đợi, số lượt retry và phân vị thời gian xử lý;
        \texttt{GET /readyz} báo trạng thái worker cho giám sát.
\end{itemize}

% =====================================================================
\section{Mô hình dữ liệu vector theo section}
% =====================================================================

Điểm đặc biệt của hệ thống là mỗi tài liệu \emph{không} được nhúng thành một vector duy
nhất mà tách theo section (kỹ năng, kinh nghiệm, dự án, học vấn, tóm tắt) cộng thêm một
vector cho toàn văn. Nhờ đó khi so một JD với một CV, ta so được ``yêu cầu kỹ năng''
trùng với ``phần kỹ năng'' thay vì so toàn văn mơ hồ.

\begin{figure}[htbp]
\centering
\resizebox{0.95\textwidth}{!}{%
\begin{tikzpicture}[node distance=0.55cm and 1.1cm, font=\scriptsize]
  \node[box, fill=cbluel, minimum width=3.2cm] (doc) at (0,-1.9)
        {\textbf{CV id=1234, version=1}\\(payload đã phân tích)};
  \foreach \i/\nm in {0/{skills}, 1/{experience}, 2/{projects},
                      3/{education}, 4/{summary}, 5/{\_full}} {
    \node[box, fill=cgreenl, minimum width=2.1cm] (sec\i) at (5.4,{-0.8*\i}) {\nm};
    \draw[arr] (doc.east) -- (sec\i.west);
  }
  \node[infra, minimum height=4.6cm, minimum width=3.6cm] (col) at (10.4,-2)
        {Collection \texttt{frevia\_cvs}\\[2pt]
         \texttt{cv:1234:v1:skills}\\
         \texttt{cv:1234:v1:experience}\\
         \texttt{cv:1234:v1:projects}\\
         \texttt{cv:1234:v1:education}\\
         \texttt{cv:1234:v1:summary}\\
         \texttt{cv:1234:v1:\_full}};
  \foreach \i in {0,1,2,3,4,5} \draw[arr] (sec\i) -- ([xshift=-1.4cm]col.west |- sec\i);
  \node[store, minimum width=3.4cm] (js) at (15.2,-2.6)
        {\texttt{.documents.jsonl}\\(payload + owner + visibility)};
  \draw[arr] (col) -- node[lbl,above]{ghi cùng lúc} (js);
\end{tikzpicture}}
\caption{Mỗi CV/JD sinh ra một ``bó'' vector theo section. Vector ID có dạng
\texttt{type:id:v*n*:section}, metadata gắn kèm gồm \texttt{content\_hash},
\texttt{status=READY}, \texttt{owner\_id}, \texttt{visibility}.}
\label{fig:vector}
\end{figure}


Khi phiên bản tài liệu tăng (\texttt{version}$+1$) thì upsert sẽ tự xoá các vector cũ
cùng \texttt{document\_id} và thay bằng bộ mới; metadata \texttt{status=READY} đảm bảo
chỉ vector đã hoàn chỉnh mới được dùng cho truy hồi.

% =====================================================================
\section{Luồng ghép cặp CV -- việc làm}
% =====================================================================

Đây là phần lõi của hệ thống: cho một người dùng, tìm ra những ứng viên (hoặc tin
tuyển dụng) phù hợp nhất.

\subsection{Từ file tải lên đến kết quả ghép cặp}

Trước khi ghép cặp được, tài liệu phải đi qua một chuỗi xử lý để biến văn bản thô thành
dữ liệu máy ``hiểu'' được. Hình~\ref{fig:fullflow} mô tả toàn bộ chuỗi đó, gồm tám bước
chia làm hai giai đoạn: \textbf{xử lý tài liệu} (bước 1--7) và \textbf{ghép cặp}
(bước 8).

\begin{figure}[htbp]
\centering
\resizebox{\textwidth}{!}{%
\begin{tikzpicture}[
    font=\small,
    st/.style = {draw=cblue, line width=0.9pt, rounded corners=5pt, fill=cbluel,
                 text width=2.9cm, minimum height=1.9cm, align=center, inner sep=4pt},
    ll/.style = {draw=cpurple, line width=0.9pt, rounded corners=5pt, fill=cpurplel,
                 text width=2.9cm, minimum height=1.9cm, align=center, inner sep=4pt},
    db/.style = {draw=corange, line width=0.9pt, rounded corners=5pt, fill=corangel,
                 text width=2.9cm, minimum height=1.9cm, align=center, inner sep=4pt},
    mm/.style = {draw=cgreen, line width=1.0pt, rounded corners=5pt, fill=cgreenl,
                 text width=2.9cm, minimum height=1.9cm, align=center, inner sep=4pt},
    fa/.style = {-{Stealth[length=2.6mm]}, thick},
    fl/.style = {font=\scriptsize, align=center, inner sep=2pt, text=cgray},
  ]

  % ================== HANG 1: xu ly tai lieu ==================
  \node[st] (up) at (0,0)
        {\textbf{1. Tải lên}\\người dùng gửi\\CV / tin tuyển dụng\\(PDF, DOCX, chữ)};
  \node[st] (ex) at (4.0,0)
        {\textbf{2. Trích xuất chữ}\\đọc nội dung file:\\PDF, DOCX,\\ảnh scan thì OCR};
  \node[ll] (pa) at (8.0,0)
        {\textbf{3. Phân tích bằng AI}\\máy đọc chữ, rút ra\\kỹ năng, kinh nghiệm,\\học vấn, tóm tắt\dots};
  \node[ll] (no) at (12.0,0)
        {\textbf{4. Chuẩn hoá}\\sửa ngày tháng,\\bỏ trùng, kiểm tra\\dẫn chứng};

  % ================== HANG 2: index + matching ==================
  \node[st] (se) at (0,-4.4)
        {\textbf{5. Tách thành phần}\\kỹ năng / kinh nghiệm\\/ học vấn / tóm tắt\\+ toàn văn};
  \node[ll] (em) at (4.0,-4.4)
        {\textbf{6. Chuyển thành vector}\\mỗi phần $\to$ một dãy số\\đại diện cho ý nghĩa\\(BGE-M3)};
  \node[db] (ch) at (8.0,-4.4)
        {\textbf{7. Lưu vào kho vector}\\một tài liệu $\to$\\nhiều vector, kèm\\trạng thái \& chủ sở hữu};
  \node[mm] (rc) at (12.0,-4.4)
        {\textbf{8. Ghép cặp}\\so từng phần,\\cộng điểm, xếp hạng\\$\to$ trả kết quả};

  % ================== canh noi ==================
  \draw[fa] (up) -- (ex);
  \draw[fa] (ex) -- (pa);
  \draw[fa] (pa) -- (no);
  \draw[fa] (no.south) -- ++(0,-0.6) -- (-2.0,-2.2) -- (-2.0,-3.8) -- (up.south |- 0,-3.8) -- (se.west);
  \draw[fa] (se) -- (em);
  \draw[fa] (em) -- (ch);
  \draw[fa] (ch) -- (rc);

  % ================== nhan hai giai doan ==================
  \node[fl, text=cblue] at (6.0,1.15)
        {\textbf{Giai đoạn 1 --- xử lý tài liệu (chỉ làm một lần, khi tài liệu mới được tải lên)}};
  \node[fl, text=cgreen] at (6.0,-3.25)
        {\textbf{Giai đoạn 2 --- ghép cặp (làm mỗi khi cần tìm ứng viên)}};

  % ================== ghi chu duoi ==================
  \node[fl] at (6.0,-6.0)
        {Bước 1--7 biến văn bản thô thành các vector đã lưu sẵn;
         bước 8 dùng lại các vector đó để tìm ứng viên phù hợp.};
\end{tikzpicture}}
\caption{Luồng đầy đủ từ khi tài liệu được tải lên đến khi ghép cặp. Giai đoạn~1
(bước 1--7) xử lý và lưu trữ tài liệu một lần; giai đoạn~2 (bước 8) dùng lại dữ liệu đã
lưu để ghép cặp và trả kết quả.}
\label{fig:fullflow}
\end{figure}


Tám bước được giải thích như sau:

\begin{enumerate}[leftmargin=1.6em]
  \item \textbf{Tải lên.} Người dùng gửi một CV hoặc một tin tuyển dụng. Hệ thống nhận
        file PDF, DOCX (hoặc chữ dán trực tiếp) và lưu tạm lại để xử lý về sau.
  \item \textbf{Trích xuất chữ.} Hệ thống đọc nội dung chữ ra khỏi file: PDF đọc bằng
        thư viện \texttt{pypdf}, DOCX đọc bằng \texttt{python-docx}, file scan (ảnh) thì
        đọc bằng OCR. Sau đó dọn dẹp ký tự rác, gộp những dòng bị ngắt giữa câu.
  \item \textbf{Phân tích bằng AI.} Một mô hình ngôn ngữ chạy \emph{nội bộ} (không gửi
        dữ liệu ra ngoài) đọc đoạn chữ và rút ra thông tin có cấu trúc: kỹ năng, kinh
        nghiệm, học vấn, dự án, tóm tắt (với CV) hoặc kỹ năng yêu cầu, trách nhiệm, ngân
        sách (với tin tuyển dụng).
  \item \textbf{Chuẩn hoá.} Máy sửa lại cho thống nhất: ngày tháng về đúng một dạng, bỏ
        các mục trùng lặp, và \emph{kiểm tra dẫn chứng} --- chỉ giữ những thông tin thật
        sự xuất hiện trong văn bản gốc.
  \item \textbf{Tách thành từng phần.} Thay vì coi cả tài liệu là một khối, hệ thống
        chia thành các \emph{phần} riêng: kỹ năng, kinh nghiệm, học vấn, tóm tắt\dots{}
        và thêm một phần ``toàn văn'' gộp tất cả.
  \item \textbf{Chuyển thành vector.} Mỗi phần được mô hình BGE-M3 chuyển thành một dãy
        số (vector) đại diện cho \emph{ý nghĩa} của phần đó. Nhờ vậy máy so sánh được
        ``giống nhau về nghĩa'' chứ không chỉ giống mặt chữ.
  \item \textbf{Lưu vào kho vector.} Các dãy số được lưu lại, kèm theo trạng thái (đã
        xử lý xong), chủ sở hữu và mức hiển thị. Một tài liệu sinh ra \emph{nhiều} vector
        --- một cho mỗi phần.
  \item \textbf{Ghép cặp.} Khi cần tìm ứng viên, hệ thống dùng lại các vector đã lưu,
        so từng phần với nhau, cộng điểm chi tiết rồi trả về danh sách phù hợp nhất.
\end{enumerate}

\subsection{Các bước trong phần ghép cặp}

Bước 8 ở trên được mở rộng thành bốn bước nhỏ (Hình~\ref{fig:matching}), đi từ trên
xuống dưới:

\begin{enumerate}[leftmargin=1.6em]
  \item \textbf{Tìm ứng viên phù hợp.} Nếu bên gọi chưa có danh sách ứng viên, hệ thống
        tự tìm trong kho dữ liệu đã lưu, xếp theo mức giống nhau về nội dung và chỉ giữ
        những hồ sơ mà người dùng được phép xem. Nếu bên gọi \emph{đã} có danh sách sẵn
        thì bỏ qua bước này.
  \item \textbf{Chấm điểm mức độ hợp.} Hệ thống so nội dung của hai bên \emph{theo từng
        phần}: phần kỹ năng so với phần kỹ năng, phần kinh nghiệm so với phần kinh
        nghiệm\dots{} thay vì so cả văn bản một cách chung chung. Nhờ vậy điểm số sát ý
        nghĩa hơn.
  \item \textbf{Cộng điểm chi tiết.} Ngoài mức hợp về nội dung, hệ thống cộng thêm điểm
        theo các yếu tố rõ ràng: mức đáp ứng kỹ năng, số năm kinh nghiệm, ngân sách và
        sở thích (xem Mục~\ref{sec:hai-chieu} và~\ref{subsec:diem}).
  \item \textbf{Lọc và xếp hạng.} Cuối cùng hệ thống loại những hồ sơ không còn hợp lệ
        (đã xoá, không đúng người xem, chưa sẵn sàng), sắp xếp theo điểm giảm dần và trả
        về danh sách phù hợp nhất.
\end{enumerate}

Bước \textbf{giải thích bằng AI} là tuỳ chọn: khi bật, máy sẽ tự viết lý do vì sao một
ứng viên được xem là hợp, giúp người dùng tin tưởng kết quả hơn.

\begin{figure}[htbp]
\centering
\resizebox{0.94\textwidth}{!}{%
\begin{tikzpicture}[
    node distance=0.85cm and 0.9cm,
    font=\small,
    bx/.style = {draw=cblue, line width=0.9pt, rounded corners=5pt, fill=cbluel,
                   minimum width=6.6cm, minimum height=1.0cm, align=center,
                   inner sep=5pt},
    choice/.style = {draw=corange, line width=0.9pt, rounded corners=5pt, fill=corangel,
                   minimum width=6.6cm, minimum height=1.0cm, align=center,
                   inner sep=5pt},
    opt/.style = {draw=cpurple, line width=0.9pt, dashed, rounded corners=5pt, fill=cpurplel,
                   minimum width=4.4cm, minimum height=0.95cm, align=center,
                   inner sep=5pt},
    end/.style = {draw=cgreen, line width=1.0pt, rounded corners=5pt, fill=cgreenl,
                   minimum width=6.6cm, minimum height=1.0cm, align=center,
                   inner sep=5pt, font=\small\bfseries},
    fa/.style = {-{Stealth[length=2.4mm]}, thick},
    fl/.style = {font=\footnotesize, align=center, inner sep=2pt},
  ]

  % ---------- luong chinh ----------
  \node[bx] (in)   at (0,0)
        {\textbf{Bắt đầu}\\Người dùng tìm ứng viên cho một tin / tìm việc cho một CV};

  \node[choice] (has) at (0,-2.0)
        {\textbf{Đã có sẵn danh sách ứng viên?}};

  \node[bx] (find) at (-4.4,-4.2)
        {\textbf{Tìm ứng viên phù hợp}\\đọc dữ liệu đã lưu, lọc những hồ sơ người dùng được xem};

  \node[choice] (score1) at (0,-6.3)
        {\textbf{Chấm điểm mức độ hợp}};

  \node[bx] (score2) at (0,-8.3)
        {\textbf{Cộng điểm chi tiết}\\
         kỹ năng, kinh nghiệm, ngân sách, sở thích};

  \node[bx] (filter) at (0,-10.3)
        {\textbf{Lọc theo điều kiện bắt buộc}\\chỉ giữ hồ sơ còn hợp lệ và được phép xem};

  \node[end] (out) at (0,-12.3)
        {\textbf{Xếp hạng và trả kết quả}\\danh sách phù hợp nhất kèm lý do};

  % ---------- nhanh phu ----------
  \node[bx] (skip) at (4.9,-4.2)
        {\textbf{Dùng danh sách có sẵn}\\bỏ qua bước tìm kiếm};

  \node[opt] (expl) at (4.9,-10.3)
        {\textbf{Giải thích bằng AI}\\\emph{(tuỳ chọn)}};

  % ---------- canh noi ----------
  \draw[fa] (in) -- (has);
  \draw[fa] (has.west) -- node[fl,above]{Chưa có} ++(-4.4,0) |- (find.north);
  \draw[fa] (has.east) -- node[fl,above]{Đã có} ++(4.9,0) |- (skip.north);
  \draw[fa] (find.south) -- ++(0,-0.5) -- (0,-5.35) -- (score1.north);
  \draw[fa] (skip.south) -- ++(0,-0.5) -- (0,-5.35);
  \draw[fa] (score1) -- (score2);
  \draw[fa] (score2) -- (filter);
  \draw[fa] (filter) -- (out);
  \draw[fa, dashed] (out.east) -- node[fl,right]{nếu cần} (expl.south);

  % ---------- ghi chu nhanh ----------
  \node[fl, text=cgray, below=0.08cm of find]
        {tìm theo nội dung, xếp theo độ giống cao nhất};
  \node[fl, text=cgray, below=0.08cm of skip]
        {dùng điểm đã tính trước đó};
  \node[fl, text=cgray, below=0.08cm of score1]
        {so nội dung từng phần: kỹ năng với kỹ năng, kinh nghiệm với kinh nghiệm\dots};
  \node[fl, text=cgray, below=0.08cm of expl]
        {máy tự viết lý do vì sao hợp};
\end{tikzpicture}}
\caption{Luồng ghép cặp CV--việc làm. Hệ thống tìm ứng viên, chấm điểm mức độ hợp,
lọc theo điều kiện bắt buộc rồi trả về danh sách phù hợp nhất; bước giải thích bằng AI
là tuỳ chọn.}
\label{fig:matching}
\end{figure}


\subsection{Điểm mức độ hợp được tính thế nào?}
\label{subsec:diem}

Điểm cuối cùng là trung bình có trọng số của năm phần, mỗi phần nằm trong khoảng
$0$--$1$:

\begin{center}\small
\begin{tabular}{@{}lc l@{}}
\toprule
\textbf{Phần} & \textbf{Trọng số} & \textbf{Ý nghĩa} \\
\midrule
Mức hợp nội dung & $0.40$ & Hai tài liệu giống nhau đến đâu \\
Kỹ năng          & $0.30$ & Ứng viên có đủ kỹ năng tin tuyển dụng yêu cầu không \\
Kinh nghiệm      & $0.15$ & Số năm kinh nghiệm so với yêu cầu \\
Ngân sách        & $0.10$ & Khoảng lương/ngân sách hai bên có giao nhau không \\
Sở thích         & $0.05$ & Hình thức làm việc, địa điểm\dots{} có khớp không \\
\bottomrule
\end{tabular}
\end{center}

Chi tiết công thức:

\begin{align}
  c_k(q,d) &= \max_{y\in S_d}\ \cos\big(e(x_k),\,e(y)\big),
      \qquad x_k\in S_q \label{eq:sec}\\
  s_{\text{sem}}(q,d) &= \operatorname{agg}_{k}\ c_k(q,d), \qquad
      \operatorname{agg}\in\{\max,\operatorname{mean}\} \label{eq:sem}\\
  s(q,d) &= \alpha\cdot\widetilde{s}_{\text{sem}}(q,d) + (1-\alpha)\cdot\widetilde{s}_{\text{kw}}(q,d)
      \label{eq:fuse}\\
  S_{\text{hybrid}} &= \frac{\sum_i w_i\,S_i}{\sum_i w_i}, \qquad
      (w_{sem},w_{sk},w_{exp},w_{bud},w_{pref}) = (0.40,\,0.30,\,0.15,\,0.10,\,0.05)
      \label{eq:hybrid}
\end{align}

Trong đó $S_q,S_d$ là tập section của hai tài liệu, $e(\cdot)$ là vector BGE-M3, $c_k$
là điểm của từng phần, $\widetilde{\cdot}$ là chuẩn hoá min--max trên tập ứng viên, và
$\alpha$ là tham số pha trộn giữa mức hợp nội dung và mức khớp từ khóa (mặc định $0.5$).
Phép so theo phần ở Công thức~\eqref{eq:sec} là \texttt{cross} (lấy phần giống nhất);
chế độ \texttt{fixed} thay bằng việc so \emph{đúng} phần cùng tên.

\subsection{Hai chiều ghép cặp}
\label{sec:hai-chieu}

Hệ thống chạy được cả hai chiều trên cùng một mã nguồn (Hình~\ref{fig:bidirectional}):

\begin{itemize}[leftmargin=1.2em]
  \item \textbf{Tìm CV cho tin tuyển dụng} (\texttt{jd2cv}): đầu vào là một tin tuyển
        dụng, kết quả là danh sách CV phù hợp nhất.
  \item \textbf{Tìm việc cho CV} (\texttt{cv2jd}): đầu vào là một CV, kết quả là danh
        sách tin tuyển dụng phù hợp nhất.
\end{itemize}

\begin{figure}[htbp]
\centering
\resizebox{0.96\textwidth}{!}{%
\begin{tikzpicture}[
    font=\small,
    inb/.style = {draw=cblue, line width=0.9pt, rounded corners=5pt, fill=cbluel,
                  minimum width=3.0cm, minimum height=1.0cm, align=center, inner sep=5pt},
    outb/.style = {draw=cgreen, line width=0.9pt, rounded corners=5pt, fill=cgreenl,
                  minimum width=3.4cm, minimum height=1.0cm, align=center, inner sep=5pt},
    mid/.style = {draw=corange, line width=0.9pt, rounded corners=5pt, fill=corangel,
                  minimum width=4.4cm, minimum height=1.6cm, align=center, inner sep=5pt},
    fa/.style = {-{Stealth[length=2.4mm]}, thick},
    fl/.style = {font=\footnotesize, align=center, inner sep=2pt},
  ]

  % -------- chieu 1 --------
  \node[inb]  (i1) at (0,0)     {\textbf{Một tin tuyển dụng}};
  \node[mid]  (m1) at (5.4,0)   {\textbf{Hệ thống tìm \& chấm điểm}\\công thức giống nhau\\cho cả hai chiều};
  \node[outb] (o1) at (11.0,0)  {\textbf{Danh sách CV}\\\small phù hợp nhất};

  % -------- chieu 2 --------
  \node[inb]  (i2) at (0,-3.0)  {\textbf{Một CV}};
  \node[outb] (o2) at (11.0,-3.0) {\textbf{Danh sách tin tuyển dụng}\\\small phù hợp nhất};

  \draw[fa] (i1) -- node[fl,above]{\texttt{jd2cv}} (m1);
  \draw[fa] (m1) -- (o1);
  \draw[fa] (i2) -- (m1);
  \draw[fa] (m1) -- (o2);

  % -------- ghi chu --------
  \node[fl, text=cgray] at (5.4,-4.3)
        {Chỉ khác nhau ở \textbf{vai trò}: tin tuyển dụng luôn là ``bên yêu cầu'', CV luôn là ``bên đáp ứng''.};
\end{tikzpicture}}
\caption{Hai chiều ghép cặp dùng chung một cách tính điểm: một tin tuyển dụng đi tìm
CV, và một CV đi tìm tin tuyển dụng.}
\label{fig:bidirectional}
\end{figure}


Điểm giống nhau: công thức chấm điểm, cách lọc và cách giải thích \emph{hoàn toàn như
nhau}. Điểm khác nhau duy nhất là \emph{vai trò} của hai tài liệu --- tin tuyển dụng
luôn đóng vai ``bên yêu cầu'', CV luôn đóng vai ``bên đáp ứng'' --- và kho dữ liệu được
tra cứu (CV hay tin tuyển dụng). Nhờ vậy chỉ cần viết và kiểm thử một lần.



% =====================================================================
\section{Hợp đồng API (tóm tắt)}
% =====================================================================

\begin{table}[htbp]
\centering\small
\begin{tabular}{@{}llp{8.6cm}@{}}
\toprule
\textbf{Method \& Path} & \textbf{Trạng thái} & \textbf{Vai trò} \\
\midrule
\texttt{POST /api/v1/ai/analyze-cv} & 200 & Phân tích CV đồng bộ (trả JSON, không lưu) \\
\texttt{POST /internal/extract-cv}  & 200 & Trích xuất + phân tích CV (nội bộ) \\
\texttt{POST /internal/extract-job} & 200 & Trích xuất + phân tích JD (nội bộ) \\
\texttt{POST /internal/index}       & 200 & Nhúng + upsert vector cho danh sách payload \\
\texttt{POST /internal/index/delete}& 200 & Xoá vector theo id/type \\
\texttt{POST /internal/recommend}   & 200 & Xếp hạng hybrid (kèm \texttt{explain}) \\
\texttt{POST /internal/explain}     & 200 & Giải thích 1 cặp (job, cv) bằng LLM \\
\texttt{POST /internal/rerank}      & 200 & Chấm lại theo rubric 0--100 \\
\texttt{POST /internal/embeddings}  & 200 & Sinh vector thô cho danh sách chuỗi \\
\midrule
\texttt{POST /internal/jobs}        & 202 & Tạo job xử lý file (multipart) $\to$ \texttt{job\_id} \\
\texttt{GET  /internal/jobs/\{id\}}  & 200 & Trạng thái + tiến độ + cảnh báo \\
\texttt{DELETE /internal/jobs/\{id\}}& 200 & Hủy job chưa xong \\
\texttt{POST /internal/jobs/\{id\}/retry} & 200 & Đưa job \texttt{FAILED} về \texttt{PENDING} \\
\texttt{GET  /internal/jobs/stats}  & 200 & Số job theo trạng thái, hàng đợi, phân vị thời gian \\
\texttt{POST /internal/jobs/reconcile} & 200 & Gỡ job kẹt (idempotent) \\
\midrule
\texttt{GET /healthz}, \texttt{GET /readyz} & 200 & Sống/khỏe (ping Ollama, độ sâu hàng đợi) \\
\bottomrule
\end{tabular}
\caption{Các endpoint của AI service. Mọi đường \texttt{/internal/*} đều bắt buộc
\texttt{X-Internal-Api-Key} khi bật \texttt{INTERNAL\_API\_KEY\_REQUIRED=1}.}
\label{tab:api}
\end{table}

Mẫu phản hồi tiêu biểu của job (rút gọn):

\begin{verbatim}
{
  "job_id": "j_9f2c...", "status": "PROCESSING",
  "document_type": "cv", "document_id": "1234", "version": 3,
  "step": "PARSING", "progress": 55,
  "attempt_count": 1, "warnings": [],
  "duration_ms": null, "result": null
}
\end{verbatim}

% =====================================================================
\section{Cấu hình \& công nghệ}
% =====================================================================

\begin{table}[htbp]
\centering\footnotesize
\begin{tabular}{@{}llll@{}}
\toprule
\textbf{Nhóm} & \textbf{Thành phần} & \textbf{Lựa chọn} & \textbf{Lý do} \\
\midrule
Web frontend & Framework & Next.js & (repo riêng) \\
API nghiệp vụ & Backend & NestJS + Prisma & (repo riêng) \\
AI service & Ngôn ngữ & Python 3.13 + FastAPI & Async, OpenAPI tự sinh \\
LLM & Mô hình & \texttt{qwen3:4b} (Ollama) & Không mất dữ liệu ra ngoài \\
Embedding & Mô hình & \texttt{BAAI/bge-m3} (1024 chiều, CPU) & Đa ngôn ngữ, nhiều chức năng \\
Vector store & Cơ sở & ChromaDB persistent (HNSW, cosine) & Nhúng thư viện, đơn giản \\
Job store & Cơ sở & SQLite (WAL) & PG đang chỉ đọc; interface để chuyển Postgres \\
Từ khóa & Thuật toán & BM25 (\texttt{rank\_bm25}) & Baseline kinh điển, giải thích được \\
Kiểm thử & Công cụ & pytest, ruff & 138 test tự động \\
\bottomrule
\end{tabular}
\caption{Công nghệ chính của AI service.}
\label{tab:stack}
\end{table}

Các biến môi trường quan trọng của hàng đợi (mặc định trong dấu ngoặc):

\begin{table}[htbp]
\centering\small
\begin{tabular}{@{}llp{7.6cm}@{}}
\toprule
\textbf{Biến} & \textbf{Mặc định} & \textbf{Vai trò} \\
\midrule
\texttt{OLLAMA\_MODEL} & -- (bắt buộc) & Mô hình LLM dùng cho parse/explain/rerank \\
\texttt{EMBED\_MODEL} & \texttt{BAAI/bge-m3} & Mô hình sinh embedding \\
\texttt{CHROMA\_PATH} & \texttt{None} & Nơi lưu Chroma; \texttt{None} = bộ nhớ (dùng cho test) \\
\texttt{INTERNAL\_API\_KEY\_REQUIRED} & \texttt{0} & Bật = ``fail-closed'' khi chưa cấu hình key \\
\texttt{PROCESSING\_WORKER\_ENABLED} & \texttt{0} & Bật worker nền trong process FastAPI \\
\texttt{PROCESSING\_DB\_PATH} & \texttt{data/processing.db} & Đường dẫn SQLite của hàng đợi \\
\texttt{PROCESSING\_FILES\_DIR} & \texttt{data/processing\_files} & Thư mục file tạm \\
\texttt{JOB\_MAX\_ATTEMPTS} & 3 & Số lượt thử tối đa cho lỗi tạm thời \\
\texttt{JOB\_STALE\_SECONDS} & 300 & Ngưỡng coi job \texttt{PROCESSING} là kẹt \\
\texttt{JOB\_RETRY\_BASE\_DELAY\_S} & 1.0 & Hệ số backoff $2^{n}\cdot\Delta$ \\
\texttt{PROCESSING\_SKIP\_REEMBED} & 1 & Bật bỏ qua nhúng lại khi nội dung không đổi \\
\bottomrule
\end{tabular}
\caption{Cấu hình hàng đợi xử lý (xem đầy đủ tại \texttt{apps/ai/.env.example}).}
\label{tab:env}
\end{table}

% =====================================================================
\section{Kiểm thử \& đánh giá}
% =====================================================================

\begin{itemize}[leftmargin=1.2em]
  \item \textbf{Kiểm thử tự động:} 138 test pytest bao phủ: trích xuất văn bản, chuẩn hoá
        payload, xếp hạng hybrid, giải thích/rerank (LLM giả lập), máy trạng thái job,
        worker (retry, backoff, job kẹt, hủy), hợp đồng API. Chroma dùng \emph{ephemeral
        client} + \texttt{MockEmbedder} nên test chạy nhanh, không tải model thật.
  \item \textbf{Kiểm soát chất lượng mã:} \texttt{ruff} cho mã mới; tài liệu thiết kế
        (\texttt{docs/DESIGN\_processing\_queue.md}) được cập nhật đồng bộ với mã.
  \item \textbf{Đánh giá chất lượng ghép cặp:} dùng bộ chỉ số chuẩn của truy hồi
        thông tin --- Precision@k, Recall@k, nDCG@k, MRR, MAP --- trên tập có nhãn
        (xem mục sau).
\end{itemize}

% =====================================================================
\section{Đề xuất hướng nghiên cứu khoa học}
% =====================================================================

\subsection{Câu hỏi nghiên cứu}

\begin{enumerate}[leftmargin=1.5em]
  \item \textbf{RQ1 --- Giá trị của việc tách theo section:} xếp hạng theo multi-vector
        từng section có tốt hơn nhúng toàn văn (single-vector) và BM25 thuần hay không?
  \item \textbf{RQ2 --- Pha trộn semantic -- từ khóa:} giá trị $\alpha$ nào tối ưu
        (\texttt{eq.~\ref{eq:fuse}}) trên dữ liệu tiếng Việt/kép Việt -- Anh?
  \item \textbf{RQ3 --- Vai trò của LLM ở cuối pipeline:} rerank theo rubric và giải
        thích evidence-grounded có cải thiện thứ tự xếp hạng và độ tin cậy của người
        dùng không (đo bằng LLM-as-judge và khảo sát người thật)?
  \item \textbf{RQ4 --- Chi phí -- hiệu năng:} hàng đợi nền với bỏ qua nhúng lại
        (\texttt{content\_hash}) tiết kiệm bao nhiêu thời gian/token so với xử lý đồng bộ
        và xử lý lại toàn bộ?
\end{enumerate}

\subsection{Ma trận thí nghiệm đề xuất}

\begin{table}[htbp]
\centering\small
\begin{tabular}{@{}llll@{}}
\toprule
\textbf{Nhóm biến} & \textbf{Tùy chọn} & \textbf{Chỉ số} \\
\midrule
Phương pháp & BM25; single-vector; section-aware (mean/max); chunk max-sim & nDCG@10, MRR \\
Tham số & $\alpha\in\{0,0.3,0.5,0.7,1.0\}$; pool $P\in\{50,100,200\}$ & nDCG@10, độ trễ \\
Embedding & bge-m3; multilingual-e5-small; MiniLM đa ngôn ngữ & nDCG@10, thời gian nhúng \\
Tầng LLM & không rerank; rerank; rerank + explain & MRR, điểm LLM-as-judge (1--5) \\
Hàng đợi & đồng bộ; nền toàn bộ; nền + skip re-embed & thời gian phản hồi, token/ job \\
\bottomrule
\end{tabular}
\caption{Ma trận thí nghiệm (mỗi cấu hình chạy 3 lần, lấy trung bình).}
\label{tab:exp}
\end{table}

\subsection{Dữ liệu \& quy trình}

\begin{itemize}[leftmargin=1.2em]
  \item Bộ dữ liệu công khai (Kaggle: \emph{Resume Dataset}, \emph{Job Descriptions})
        và dữ liệu nội bộ đã có sẵn pipeline đồng bộ từ PostgreSQL sang Chroma;
        yêu cầu tối thiểu 500 cặp CV--JD có nhãn ``phù hợp / không phù hợp''.
  \item Chia $70/10/20$ (train/dev/test); tham số $\alpha$, pool, ngưỡng lọc ràng buộc
        được tinh chỉnh trên dev, chỉ chốt một lần trên test.
  \item Đánh giá con người: 3--5 người đánh giá top-10 gợi ý theo thang 1--5 cho
        \emph{mức độ phù hợp} và \emph{mức độ hữu ích của phần giải thích}.
\end{itemize}

\subsection{Lộ trình 12 tuần (đề xuất)}

\begin{table}[htbp]
\centering\small
\begin{tabular}{@{}clp{9.4cm}@{}}
\toprule
\textbf{Tuần} & \textbf{Giai đoạn} & \textbf{Nội dung} \\
\midrule
1--2   & Chuẩn bị   & Đọc tài liệu, dựng bộ dữ liệu có nhãn, hoàn tất baseline BM25 \\
3--5   & Thí nghiệm & Chạy section-aware vs single-vector (RQ1), tinh chỉnh $\alpha$ (RQ2) \\
6--8   & LLM        & Thí nghiệm rerank + explain, đánh giá LLM-as-judge (RQ3) \\
9--10  & Hệ thống   & Đo chi phí hàng đợi vs đồng bộ, đo độ trễ đầu-cuối (RQ4) \\
11--12 & Viết       & Viết báo cáo, vẽ biểu đồ, chuẩn bị slide thuyết trình \\
\bottomrule
\end{tabular}
\caption{Lộ trình nghiên cứu đề xuất.}
\label{tab:timeline}
\end{table}

\subsection{Rủi ro \& biện pháp}

\begin{table}[htbp]
\centering\small
\begin{tabular}{@{}lp{9.6cm}@{}}
\toprule
\textbf{Rủi ro} & \textbf{Biện pháp} \\
\midrule
Thiếu dữ liệu có nhãn & Dùng dữ liệu công khai + gán nhãn thủ công theo mẫu rubric; dùng
                       LLM-assisted labeling rồi kiểm tra chéo 20\% \\
LLM chấm không nhất quán & Chạy 3 lần, tính đồng thuận (kappa); chốt prompt cố định \\
Ollama/BGE-M3 chậm trên CPU & Giữ worker 1 thread, đo riêng từng bước, dùng pool nhỏ; nếu
                       cần chuyển sang GPU/API embedding \\
Kết quả không có ý nghĩa thống kê & Kiểm định Wilcoxon signed-rank giữa các cấu hình \\
\bottomrule
\end{tabular}
\caption{Rủi ro dự kiến.}
\label{tab:risk}
\end{table}

\subsection{Tài liệu tham khảo chính}

\begin{itemize}[leftmargin=1.2em, itemsep=2pt]
  \item Chen, J. et al. (2024). \emph{BGE M3-Embedding: Multi-Lingual, Multi-Functionality,
        Multi-Granularity Text Embeddings Through Self-Knowledge Distillation}. ACL Findings.
  \item Robertson, S. \& Zaragoza, H. (2009). \emph{The Probabilistic Relevance Framework:
        BM25 and Beyond}. Foundations and Trends in IR.
  \item Karpukhin, V. et al. (2020). \emph{Dense Passage Retrieval for Open-Domain Question
        Answering}. EMNLP.
  \item Nogueira, R. \& Cho, K. (2019). \emph{Passage Re-ranking with BERT}. arXiv:1901.04085.
  \item Zheng, L. et al. (2023). \emph{Judging LLM-as-a-Judge with MT-Bench and Chatbot Arena}.
        NeurIPS Datasets and Benchmarks.
  \item Tài liệu thiết kế nội bộ: \texttt{docs/DESIGN\_processing\_queue.md},
        \texttt{docs/WORKFLOW.md}, \texttt{plan.md} (trong repo).
\end{itemize}

% =====================================================================
\section*{Phụ lục: cấu trúc thư mục liên quan}
% =====================================================================
\addcontentsline{toc}{section}{Phụ lục: cấu trúc thư mục liên quan}

\begin{verbatim}
section-aware-cv-job-matching/
├── apps/ai/                     # AI service (FastAPI)
│   ├── app/
│   │   ├── main.py              # lifespan, worker thread, /healthz, /readyz
│   │   ├── routes.py            # 17 endpoint (analyze-cv, /internal/*)
│   │   ├── schemas.py           # Pydantic models
│   │   ├── security.py          # API key + fail-closed
│   │   ├── text.py              # PDF/DOCX/OCR, magic bytes, làm sạch
│   │   ├── llm.py               # client Ollama (retry, timeout)
│   │   ├── embeddings.py        # BGE-M3 (+ MockEmbedder cho test)
│   │   ├── normalize.py         # chuẩn hoá CV/JD -> payload
│   │   ├── recommend.py         # xếp hạng hybrid, index, explain
│   │   ├── explain.py           # giải thích + rerank bằng LLM
│   │   ├── index_store.py       # payload JSONL cạnh Chroma
│   │   └── processing/          # hang doi nen
│   │       ├── store.py         # SQLite + may trang thai job
│   │       ├── pipeline.py      # 4 buoc EXTRACTING -> FINALIZING
│   │       └── worker.py        # thread nen, retry, backoff, stale
│   ├── tests/                   # 135 test pytest
│   └── .env.example
├── src/cvjd/                    # thu vien loi (chay duoc tu CLI)
│   ├── schema.py, io.py
│   ├── data/                    # payload, preprocessing
│   ├── models/                  # bm25, chunking, embedding, matching
│   ├── ranking/                 # engine, hybrid, structured, constraints
│   ├── retrieval/               # retriever
│   └── storage/                 # chroma, postgres, outbox
└── docs/
    ├── DESIGN_processing_queue.md
    ├── WORKFLOW.md
    ├── kien_truc_he_thong.tex   # tai lieu nay (chi \input so do)
    └── so_do/                   # 10 so do TikZ tach rieng
        ├── fig01_tong_the.tex … fig07_vector_theo_section.tex
        ├── fig08a_luong_day_du.tex       # luong day du: upload -> tra ket qua
        ├── fig08_luong_ghep_cap.tex      # 4 buoc trong phan ghep cap
        ├── fig09_hai_chieu.tex           # jd2cv vs cv2jd
        ├── build_png.py         # sinh PNG cho tung so do
        ├── png/                 # anh PNG da sinh (200 dpi)
        └── _build/              # PDF trung gian (xoa duoc)
\end{verbatim}

\end{document}
